\chapter{概率计算}

\section{基础概率与条件概率}
\noindent\textbf{要点与方法.}\enspace
条件独立与无条件独立要分开判断。潜在硬币类型、共同故障或筛选事件都会改变观测的联合分布。按实际记录写出各情形的似然，再用全概率与贝叶斯公式归一化；若任务是选择行动，还须比较包含错误损失的条件期望收益。

对称性可简化事件概率，也可构造相关样本来降低估计方差，但零点原子不能遗漏。停止抽样时应使用相应的停止次数分布检验无偏性。对联合高斯量，正交变换可分离筛选方向；筛选后的相关性则须重新计算条件矩。
\par\smallskip



\problem{P02}{对称分布的反向配对估计}{对称性与方差}
随机变量 $S$ 与 $-S$ 同分布，记 $\tau=\Prob(S=0)$。为了估计 $p=\Prob(S>0)$，比较两个估计量：用两个独立样本 $S_1,S_2$ 构造
\[
A=\frac{\ind\{S_1>0\}+\ind\{S_2>0\}}2,
\qquad B=\frac{\ind\{S>0\}+\ind\{-S>0\}}2.
\]
证明二者无偏，求其方差，并说明何时反向配对完全消除随机误差。
\solution 对称性给出 $\Prob(S>0)=\Prob(S<0)=(1-\tau)/2$，所以 $\E A=\E B=p$。$A$ 中的两个指示变量独立，故
\[
\Var A=\frac{p(1-p)}2=\frac{1-\tau^2}{8}.
\]
而 $S$ 与 $-S$ 不可能同时为正；除 $S=0$ 外恰有一个为正。因此
\[
B=\frac{1-\ind\{S=0\}}2,
\qquad \Var B=\frac{\tau(1-\tau)}4.
\]
相减得到
\[
\boxed{\Var A-\Var B=\frac{(1-\tau)^2}{8}\ge0.}
\]
当 $0\le\tau<1$ 时，方差之比为 $2\tau/(1+\tau)$。特别地，$\tau=0$ 时 $B\equiv1/2$，配对估计没有随机误差。$\tau=1$ 时两种估计均恒为零；此时不能再以方差之比作比较。
\insight 配对的收益来自负相关。边际分布相同并不意味着两次计算必须独立；这里把正负事件配在一起，只留下零点事件的不确定性。

\problem{P03}{停止抽样中的比例估计}{负二项分布}
独立伯努利试验的成功概率为 $p\in(0,1)$。在第 $r$ 次成功时停止，$r\ge2$，总试验次数记为 $T$。证明 $(r-1)/(T-1)$ 是 $p$ 的无偏估计，并判断自然的样本成功比例 $r/T$ 是否无偏。
\solution 事件 $T=t$ 要求前 $t-1$ 次恰有 $r-1$ 次成功，且第 $t$ 次成功，故
\[
\Prob(T=t)=\binom{t-1}{r-1}p^r(1-p)^{t-r},\quad t\ge r.
\]
利用恒等式
\[
\frac{r-1}{t-1}\binom{t-1}{r-1}=\binom{t-2}{r-2},
\]
令 $j=t-r$，再用负二项级数，得到
\[
\begin{aligned}
\E\frac{r-1}{T-1}
&=p^r\sum_{j=0}^\infty\binom{j+r-2}{r-2}(1-p)^j\\
&=p^r p^{-(r-1)}=p.
\end{aligned}
\]
两种估计量逐点之差为
\[
\frac rt-\frac{r-1}{t-1}=\frac{t-r}{t(t-1)}\ge0.
\]
又因 $\Prob(T>r)>0$，取期望后严格大于零。因此 $\boxed{\E(r/T)>p}$，它向上有偏。若允许边界 $p=1$，则 $T=r$，两种估计均等于 $1$。
\insight 停止规则强制最后一次为成功，固定样本量下的比例估计不能直接照搬。去掉这次必然成功，再用前 $T-1$ 次构造比例，恰好消除偏差。

\problem{P30}{固定来源与逐次换源}{贝叶斯辨识}
实验使用两种硬币，正面概率分别为 $1/4$ 和 $3/4$。协议 $F$ 先等概率选定一种，此后一直使用它；协议 $R$ 则每次抛掷前独立等概率重选一种。给定所用硬币，各次抛掷独立。两种协议的先验概率各为 $1/2$。分别在观察到有序结果 $HH$ 与 $HT$ 后，求使用协议 $F$ 的后验概率。再求两种协议下前两次正面指示变量的协方差，解释一次观察为何无法区分协议。
\solution 两种协议的每次边际正面概率都是 $1/2$，所以单次结果的似然相同，后验仍为 $1/2$。连续两次的似然却不同：
\[
\begin{array}{c|cc}
&HH&HT\\ \hline
F&\tfrac12[(1/4)^2+(3/4)^2]=5/16&3/16\\
R&1/4&1/4
\end{array}
\]
先验相等时，后验等于本协议似然除以两协议似然之和。因此
\[
\boxed{\Prob(F\mid HH)=\frac59,\qquad
\Prob(F\mid HT)=\frac37.}
\]
令 $I_1,I_2$ 为正面指示变量。在 $F$ 下，潜在正面概率 $P$ 在 $1/4,3/4$ 间等概率取值。条件独立给出
\[
\Cov(I_1,I_2)=\E(P^2)-(\E P)^2=\Var(P)=\frac1{16}.
\]
在 $R$ 下，两次重新选择及抛掷均独立，故协方差为零。
\insight 两种机制可以具有完全相同的单次分布，却有不同的序列依赖性。连续相同结果支持固定来源，交替结果相对支持逐次换源。

\problem{P31}{允许同时发生的竞争事件}{联合分布分解}
每轮中，信号 $A$ 以概率 $\alpha$ 出现，信号 $B$ 以概率 $\beta$ 出现，$0<\alpha,\beta<1$；同轮的两信号独立，各轮也独立。在首次至少出现一个信号时停止，轮数记为 $T$，停止原因 $J$ 分为“仅 $A$”“仅 $B$”“同时出现”。求 $(T,J)$ 的联合分布，判断 $T$ 与 $J$ 是否独立，并求 $T$ 的均值与方差。
\solution 记 $q=(1-\alpha)(1-\beta)$、$s=1-q=\alpha+\beta-\alpha\beta$。三种停止原因的单轮概率分别为
\[
w_A=\alpha(1-\beta),\quad w_B=(1-\alpha)\beta,
\quad w_{AB}=\alpha\beta,
\]
其和为 $s$。停止于第 $n$ 轮要求此前均无信号，故
\[
\Prob(T=n,J=j)=q^{n-1}w_j
=(s q^{n-1})\frac{w_j}{s},\quad n\ge1.
\]
右侧分别是 $T$ 与 $J$ 的边际概率，因此二者独立。特别地，已知停止时两个信号同时出现，并不改变等待时间的分布。对几何级数求导可得
\[
\boxed{\E T=\frac1s,\qquad \Var T=\frac q{s^2},
\qquad \Prob(J=j)=\frac{w_j}{s}.}
\]
\pitfall 同时出现的概率为 $\alpha\beta$，不能把总停止概率写成 $\alpha+\beta$。独立性也不是因为“时间与原因不同”，而是来自联合概率的乘积分解。

\problem{P33}{共同故障下的多数表决}{混合模型}
二元真值 $T$ 满足 $\Prob(T=1)=\pi\in(0,1)$。三个传感器以概率 $c\in[0,1]$ 遇到共同故障，此时全部输出相反真值；无共同故障时，三者独立以概率 $\varepsilon\in(0,1/2)$ 翻转真值。故障机制与真值独立。求多数表决的错误率、观察到多数输出为 $1$ 后的真值后验，并与单个传感器比较。另求任意两个传感器错误指示变量的协方差。
\solution 无共同故障时，多数出错要求恰有两个或三个传感器出错，其概率为
\[
3\varepsilon^2(1-\varepsilon)+\varepsilon^3
=3\varepsilon^2-2\varepsilon^3.
\]
因此多数错误率与单个错误率分别为
\[
e=c+(1-c)(3\varepsilon^2-2\varepsilon^3),\quad
e_1=c+(1-c)\varepsilon.
\]
两者之差
\[
e_1-e=(1-c)\varepsilon(1-\varepsilon)(1-2\varepsilon)
\]
在 $c<1$ 时为正。共同故障部分不能通过增加表决者消除。设 $M$ 为多数输出，由两类真值下相同的翻转率，
\[
\boxed{\Prob(T=1\mid M=1)=
\frac{\pi(1-e)}{\pi(1-e)+(1-\pi)e}.}
\]
对错误指示变量 $E_i,E_j$，联合出错概率为 $c+(1-c)\varepsilon^2$，故
\[
\boxed{\Cov(E_i,E_j)=c+(1-c)\varepsilon^2-e_1^2
=c(1-c)(1-\varepsilon)^2.}
\]
\insight 先按共同故障分层，再在无故障层内使用二项分布。把三个传感器直接当作独立变量，会漏掉无法由多数表决平均掉的风险。

\problem{P12}{总和筛选引起的相关性}{正交变换与截断矩}
$X,Z$ 为独立标准正态变量。对固定 $c\ge0$，只保留满足 $X+Z>c$ 的样本。求保留样本中 $X,Z$ 的均值、方差、协方差和相关系数，证明相关系数为负，并算出 $c=0$ 的结果。记标准正态密度和分布函数为 $\phi,\Phi$。
\solution 作正交变换
\[
S=\frac{X+Z}{\sqrt2},\quad D=\frac{X-Z}{\sqrt2},
\quad a=\frac c{\sqrt2}.
\]
$S,D$ 仍是独立标准正态。筛选只限制 $S>a$，所以 $D$ 的分布及其与 $S$ 的独立性均保留。由 $\phi'(s)=-s\phi(s)$ 分部积分，
\[
\lambda:=\E(S\mid S>a)=\frac{\phi(a)}{1-\Phi(a)},\quad
\E(S^2\mid S>a)=1+a\lambda.
\]
记 $v=1+a\lambda-\lambda^2$。由 $X=(S+D)/\sqrt2$、$Z=(S-D)/\sqrt2$，条件矩为
\[
\boxed{\begin{gathered}
\E X=\E Z=\lambda/\sqrt2,\qquad
\Var X=\Var Z=(v+1)/2,\\
\Cov(X,Z)=(v-1)/2,\qquad
\operatorname{Corr}(X,Z)=\frac{v-1}{v+1}.
\end{gathered}}
\]
此处所有矩均针对筛选后的总体。由于条件下 $S>a$，有 $\lambda>a$；又 $\lambda>0$，故 $v=1-\lambda(\lambda-a)<1$，相关系数严格为负。截断分布非退化，亦有 $v>0$。

$c=0$ 时 $\lambda=\sqrt{2/\pi}$、$v=1-2/\pi$，于是均值为 $1/\sqrt\pi$、方差为 $1-1/\pi$、协方差为 $-1/\pi$、相关系数为 $-1/(\pi-1)$。
\insight 筛选依赖两个变量的总和，破坏了原来的独立性。沿总和与差值两个方向分解，可以把二维条件计算化为一个截断正态的矩。

\problem{P13}{允许放弃的贝叶斯预测}{后验决策}
硬币的正面概率 $p$ 先验均匀分布于 $(0,1)$；给定 $p$ 后，各次抛掷独立。观察 $k\ge0$ 次抛掷中的正面数 $h$ 后，可以预测下一次为正面、预测反面，或放弃。预测正确收益 $1$，错误损失 $\lambda>1$，放弃收益 $0$。求最优规则及其观察前的期望收益；计算 $k=4,\lambda=2$ 的情形。
\solution 后验密度与 $p^h(1-p)^{k-h}$ 成正比，因此下一次正面概率为
\[
r=\frac{\int_0^1p^{h+1}(1-p)^{k-h}\dd p}
{\int_0^1p^h(1-p)^{k-h}\dd p}
=\frac{h+1}{k+2}.
\]
预测正面、反面、放弃的条件收益依次为
\[
(\lambda+1)r-\lambda,\qquad 1-(\lambda+1)r,\qquad0.
\]
故 $r>\lambda/(\lambda+1)$ 时预测正面，$r<1/(\lambda+1)$ 时预测反面，其余情形放弃；在端点处相应预测与放弃无差别。

先验预测分布满足
\[
\Prob(H=h)=\binom kh\int_0^1p^h(1-p)^{k-h}\dd p
=\frac1{k+1},\quad 0\le h\le k.
\]
由 $h$ 与 $k-h$ 的对称性，令 $[z]_+=\max(z,0)$，最优期望收益为
\[
\boxed{V_k=\frac{2}{(k+1)(k+2)}
\sum_{h=0}^k[(\lambda+1)(h+1)-\lambda(k+2)]_+.}
\]
$k=4,\lambda=2$ 时，仅 $h=0,4$ 严格值得预测，条件收益均为 $1/2$；其他观测选择放弃（$h=1,3$ 处可作零收益预测），故 $V_4=1/5$。
\insight 先计算后验预测概率，再比较各动作的收益。非对称损失会产生一段放弃区间，不能只选概率较大的一面。


\par\Needspace{8\baselineskip}
\subsection*{小节练习}
\exercise{OP02}{两组硬币的正面数比较}
甲独立抛掷公平硬币 $n$ 次，乙独立抛掷公平硬币 $2n+1$ 次，其中 $n\ge1$；两人的所有抛掷相互独立。记甲、乙的正面次数分别为 $X,Y$，求 $\Prob(Y>2X)$。

\exercise{OP30}{连续两次正面后的硬币类型}
两枚硬币中，一枚公平硬币有一面正面、一面反面，另一枚两面都是正面。等概率选取一枚，随后独立抛掷选中的同一枚硬币两次。已知两次都是正面，求选中双正面硬币的条件概率。


\section{计数与期望}
\noindent\textbf{要点与方法.}\enspace
求完整分布时，数清等可能的排列、子集或间隔配置；只求均值时，常可把总数拆成指示变量之和，无须独立。方差还需两事件的联合概率，可按对象是否重叠分类。圆环计数要说明标记方法及每个对象被重复计数的次数。

交换性不等于独立；泊松总次数带来的独立分解也须由联合分布证明。项数依赖观测时，应保留这种依赖，用首步分解求随机和的矩。

高维体积可写成独立均匀样本满足约束的概率。大数定律定位阈值；在均值抵消且方差有限非零时，再用中心极限定理分析阈值附近的波动。
\par\smallskip



\problem{P05}{第一次重复抽样}{生日碰撞}
从 $m\ge1$ 种等可能类型中独立重复抽样。令 $T$ 为第一次抽到已出现类型的时刻，$L$ 为该次重复的类型。求 $T$ 的分布与期望，并判断 $L$ 是否与 $T$ 独立。停止时一共见过多少种不同类型？
\solution 记下降阶乘 $(m)_k=m(m-1)\cdots(m-k+1)$，约定 $(m)_0=1$。$T>k$ 等价于前 $k$ 次全部不同，故
\[
\Prob(T>k)=\frac{(m)_k}{m^k},\quad0\le k\le m,
\]
而 $k\ge m+1$ 时概率为零。先有 $t-1$ 个不同类型，再重复其中之一，得到
\[
\boxed{\Prob(T=t)=\frac{(m)_{t-1}}{m^{t-1}}\frac{t-1}{m},
\quad2\le t\le m+1.}
\]
利用非负整数变量的尾和公式，
\[
\boxed{\E T=\sum_{k=0}^m\frac{(m)_k}{m^k}.}
\]
对于固定 $t$，重命名类型不改变任何抽样概率，故 $\Prob(T=t,L=\ell)$ 对 $\ell=1,\ldots,m$ 全部相等，其和为 $\Prob(T=t)$。于是
\[
\Prob(T=t,L=\ell)=\Prob(T=t)/m.
\]
所以 $L$ 均匀分布且与 $T$ 独立。最后一次没有增加新类型，因此停止时的不同类型数恰为 $T-1$，其期望为 $\E T-1$。$m=1$ 时 $T=2$，公式仍成立。
\insight 第一次重复是“尚未发生碰撞”的补事件。先数无重复序列，比直接枚举重复类型及其历史位置更简洁。

\problem{P17}{随机指派中的固定点}{容斥与指示变量}
$n\ge2$ 个任务均匀随机指派给 $n$ 个人，每人恰得一个任务。任务 $i$ 原本属于第 $i$ 人。令 $F$ 为获得自己任务的人数。求 $F$ 的完整分布、均值和方差，并解释为何 $F=n-1$ 不可能发生。
\solution 指派对应于均匀随机排列。若恰有 $k$ 个固定点，先选出这些位置，剩下 $n-k$ 个位置必须构成错排。对 $r$ 个位置，容斥给出错排数
\[
D_r=\sum_{j=0}^r(-1)^j\binom rj(r-j)!
=r!\sum_{j=0}^r\frac{(-1)^j}{j!},\quad D_0=1.
\]
故
\[
\boxed{\Prob(F=k)=\frac{\binom nkD_{n-k}}{n!}
=\frac1{k!}\sum_{j=0}^{n-k}\frac{(-1)^j}{j!},
\quad0\le k\le n.}
\]
设 $I_i$ 表示第 $i$ 人获得自己的任务。固定一个位置后仍有 $(n-1)!$ 个排列，固定两个位置后有 $(n-2)!$ 个，所以
\[
\E I_i=\frac1n,\qquad \E(I_iI_j)=\frac1{n(n-1)}\quad(i\ne j).
\]
于是 $\E F=1$，且
\[
\E F^2=\sum_i\E I_i+2\sum_{i<j}\E(I_iI_j)=2,
\qquad \boxed{\Var F=1.}
\]
若已有 $n-1$ 人拿到自己的任务，最后一人也只能拿到自己的任务。相应地，分布公式中的 $D_1=0$。
\insight 完整分布需要容斥；均值和方差只需一、二阶联合概率，不必先对完整分布求和。

\problem{P18}{随机排列的短循环}{组合计数与阶乘矩}
在 $n$ 个元素的均匀随机排列中，记 $C_k$ 为长度恰为 $k$ 的循环数，$1\le k\le n$。求 $\E C_k$、$\Var C_k$，以及 $j\ne k$ 时的 $\Cov(C_j,C_k)$。说明长度受到总元素数限制时，为什么协方差会改变。
\solution 对每个 $k$ 元子集 $A$，令 $I_A$ 表示这些元素恰好构成一个循环。其内部循环有 $(k-1)!$ 种，剩余元素任意排列，所以
\[
\Prob(I_A=1)=\frac{(k-1)!(n-k)!}{n!},
\qquad \boxed{\E C_k=\frac1k.}
\]
$C_k(C_k-1)$ 计数有序的两个不同 $k$ 循环。它们必须不相交。若 $2k\le n$，其期望为
\[
\binom nk\binom{n-k}k
\frac{[(k-1)!]^2(n-2k)!}{n!}=\frac1{k^2};
\]
若 $2k>n$，两个这样的循环不能共存，阶乘矩为零。因此
\[
\boxed{\Var C_k=\frac1k-\frac{\ind\{2k>n\}}{k^2}.}
\]
类似地，$j\ne k$ 时，若 $j+k\le n$，选取两个不相交子集并分别组成循环，得到 $\E(C_jC_k)=1/(jk)$；若 $j+k>n$，乘积恒为零。于是
\[
\boxed{\Cov(C_j,C_k)=-\frac{\ind\{j+k>n\}}{jk},\quad j\ne k.}
\]
\pitfall 协方差为零不等于独立。例如所有循环长度还必须满足 $\sum_k kC_k=n$。阶乘矩把“不同循环”的计数要求直接写入了随机变量。

\problem{P19}{跳过前段后的最佳选择}{记录与最优停止}
$n\ge2$ 名候选人按均匀随机顺序到达，水平互不相同，只能观察当前候选人在已出现者中的相对排名，拒绝后不能召回。采用规则：先拒绝前 $r$ 人，再选择第一个超过此前所有人的候选人；若此后无人创新高，则选择失败。$0\le r\le n-1$，其中 $r=0$ 表示直接选第一人。求选中全体最佳者的概率，并给出最优 $r$ 的确定方法及大样本极限。
\solution $r=0$ 时成功率为 $1/n$。对 $r\ge1$，若最佳者位于 $k>r$，要一直等到他，前 $k-1$ 人中的最佳者必须在被跳过的前 $r$ 人中。其位置均匀，所以条件成功概率为 $r/(k-1)$。因此
\[
\boxed{P_r=\frac rn\sum_{j=r}^{n-1}\frac1j,\quad1\le r\le n-1.}
\]
相邻两值的差为
\[
P_{r+1}-P_r=\frac1n\left(\sum_{j=r+1}^{n-1}\frac1j-1\right),
\quad1\le r\le n-2.
\]
括号随 $r$ 严格递减，故成功率先升后降；逐次增加 $r$，直到该尾部调和和不超过 $1$ 即可确定峰值，等于 $1$ 时相邻两个值并列。还须比较 $r=0$；因 $P_1=H_{n-1}/n\ge1/n$，它不会严格优于 $r=1$。

若 $r/n\to x\in(0,1)$，调和和趋于 $-\log x$，故 $P_r\to-x\log x$。该函数在 $x=e^{-1}$ 唯一取最大值 $e^{-1}$。结合上述离散差分条件，得到最优跳过比例趋于 $1/e$，最优成功率也趋于 $1/e$。
\insight 条件化于最佳者到达时刻，把整段选择历史压缩为此前最高者的位置。优化阶段只需考察一个单调的相邻差分。

\problem{P25}{直径与偏心弦的分区}{交点计数}
单位圆内画出 $n\ge1$ 条直径，其方向独立且均匀分布于 $[0,\pi)$。再画一条预先固定的弦，它所在直线到圆心的距离为 $d\in(0,1)$。求最终分区数 $R$ 的分布、均值和方差。
\solution 以概率一，直径方向互不相同；它们全部经过圆心，形成 $2n$ 个区域。令固定弦所在直线为 $x=d$。方向角为 $\theta$ 的直径与它相交于圆内，当且仅当
\[
\frac d{|\cos\theta|}<1.
\]
因此每条直径与弦相交的概率为 $p=2\arccos(d)/\pi$。各方向独立，交点数 $K$ 服从 $\operatorname{Bin}(n,p)$。因 $d>0$，弦不经过圆心，不同直径在弦上的交点互不相同，端点相交也为零概率事件。

$K$ 个交点把新弦分成 $K+1$ 段，每段将所在区域一分为二，所以
\[
\boxed{R=2n+K+1.}
\]
于是对 $k=0,\ldots,n$，
\[
\Prob(R=2n+k+1)=\binom nkp^k(1-p)^{n-k},
\]
且
\[
\boxed{\E R=2n+1+np,\qquad \Var R=np(1-p).}
\]
\pitfall 这里原有直径共点，不能使用“每对直线贡献一个不同交点”的通用位置公式。$d=0$ 也被排除，因为此时新弦经过同一个中心，交点不再互异。

\problem{P26}{目标卡覆盖区间的长度}{位置计数与间隔}
将 $N$ 张不同卡片均匀随机排列，其中 $K$ 张为目标卡，$2\le K\le N$。记目标卡最早、最晚出现的位置为 $A,B$，令 $R=B-A$。求 $R$ 的分布与期望。
\solution $K$ 个目标位置在全部 $\binom NK$ 个 $K$ 元子集中均匀分布。给定跨度 $r$，左端点可取 $1,\ldots,N-r$；两端固定为目标，另外 $K-2$ 个目标位置须从内部 $r-1$ 个位置中选择。因此
\[
\boxed{\Prob(R=r)=\frac{(N-r)\binom{r-1}{K-2}}{\binom NK},
\quad K-1\le r\le N-1.}
\]
求均值可绕开这组有限和。把 $N-K$ 张普通卡按相邻目标卡划成 $K+1$ 个间隔，分别记为 $G_0,\ldots,G_K$，包括首尾两个外部间隔。每个非负整数解 $\sum G_i=N-K$ 唯一确定一种目标位置集合，故这些间隔可交换，
\[
\E G_i=\frac{N-K}{K+1}.
\]
又 $A=G_0+1$、$B=N-G_K$，所以
\[
\boxed{\E R=N-1-2\frac{N-K}{K+1}
=\frac{(K-1)(N+1)}{K+1}.}
\]
当 $K=N$ 时跨度恒为 $N-1$；公式同时覆盖这一边界。
\insight 分布按两个极端位置直接计数，均值按首尾空隙的交换性计算。不同目标量可以选择不同的最短推导路径。

\problem{P27}{孤立红卡的数量}{均匀整数分拆}
将 $m\ge2$ 张红卡和 $n\ge2$ 张黑卡均匀随机排列，令 $J$ 为长度恰好为 $1$ 的红色段数，首尾红色段也计入。求 $\E J$ 与 $\Var J$。
\solution 记 $N=m+n$。在 $n$ 张黑卡的前后形成 $n+1$ 个间隙，令 $G_i$ 为第 $i$ 个间隙的红卡数。每个满足 $G_i\ge0$、$\sum G_i=m$ 的整数向量唯一对应一个颜色序列，因而在 $\binom Nn$ 个向量中均匀分布。孤立红段恰好对应 $G_i=1$。

固定一个间隙为 $1$，剩余 $m-1$ 张红卡分入 $n$ 个间隙；固定两个间隙为 $1$，则剩余 $m-2$ 张分入 $n-1$ 个间隙。因此
\[
\begin{aligned}
p&:=\Prob(G_i=1)=\frac{\binom{N-2}{n-1}}{\binom Nn}
=\frac{mn}{N(N-1)},\\
p_2&:=\Prob(G_i=G_j=1)
=\frac{\binom{N-4}{n-2}}{\binom Nn}\\
&=\frac{m(m-1)n(n-1)}{N(N-1)(N-2)(N-3)},\quad i\ne j.
\end{aligned}
\]
写 $J=\sum_{i=0}^n\ind\{G_i=1\}$，逐项求一、二阶矩得到
\[
\boxed{\begin{aligned}
\E J&=(n+1)p,\\
\Var J&=(n+1)p(1-p)+n(n+1)(p_2-p^2).
\end{aligned}}
\]
\pitfall 这里间隙长度是均匀整数分拆，彼此并不独立，也不是把每张红卡独立投入间隙。两个端部间隙与内部间隙都必须计入。

\problem{P28}{圆环排列的变色次数}{循环计数}
在标号为 $1,\ldots,N$ 的圆环位置上均匀放置 $m$ 张红卡与 $n$ 张黑卡，$m,n\ge1$、$N=m+n$。令 $C$ 为相邻位置颜色不同的次数，位置 $N$ 与位置 $1$ 也相邻。求 $C$ 的完整分布及期望。$N=2$ 时按圆周的两个相邻关系分别计数。
\solution 每个红色段后接一个黑色段，颜色变化次数必为偶数。若 $C=2k$，红、黑各有 $k$ 段，故 $1\le k\le\min(m,n)$。

先标记一个红色段的起点。有 $N$ 种标记位置，再把 $m$ 分成 $k$ 个正整数作为依次出现的红段长度，把 $n$ 分成 $k$ 个正整数作为黑段长度。这样的带标记序列数为
\[
N\binom{m-1}{k-1}\binom{n-1}{k-1}.
\]
每个未标记颜色序列恰好被计数 $k$ 次，所有颜色序列共有 $\binom Nm$ 个。因此
\[
\boxed{\Prob(C=2k)=\frac{N\binom{m-1}{k-1}\binom{n-1}{k-1}}
{k\binom Nm}.}
\]
对均值，任意一条相邻关系两端异色的概率为 $2mn/[N(N-1)]$。共有 $N$ 条相邻关系，故
\[
\boxed{\E C=\frac{2mn}{N-1}.}
\]
\insight 给圆环中的一段加标记，便可把循环顺序转成线性正整数分拆；最后除以每个对象实际拥有的标记数，避免重复计数。

\problem{P29}{泊松抽样后的类别覆盖}{泊松分解}
先抽取 $N\sim\operatorname{Pois}(\lambda)$，$\lambda>0$，再独立进行 $N$ 次类别抽样，每次落入第 $i$ 类的概率为 $p_i\ge0$，$i=1,\ldots,d$，且 $\sum p_i=1$。$N$ 与抽样序列独立。令 $X_i$ 为各类出现次数，$D$ 为出现过的类别数。证明 $X_i$ 相互独立，并求 $\E D$、$\Var D$ 和 $\Cov(D,N)$。
\solution 对任意非负整数 $x_1,\ldots,x_d$，令 $s=\sum x_i$。先给定总数再用多项分布，
\[
\begin{aligned}
\Prob(X_i=x_i,\ 1\le i\le d)
&=\frac{e^{-\lambda}\lambda^s}{s!}
\frac{s!}{\prod_i x_i!}\prod_i p_i^{x_i}\\
&=\prod_{i=1}^d\frac{e^{-\lambda p_i}(\lambda p_i)^{x_i}}{x_i!}.
\end{aligned}
\]
所以各 $X_i$ 独立且服从 $\operatorname{Pois}(\lambda p_i)$；$p_i=0$ 时该计数恒为零。令 $I_i=\ind\{X_i>0\}$、$q_i=1-e^{-\lambda p_i}$，则 $D=\sum I_i$ 是独立指示变量之和，
\[
\boxed{\E D=\sum_i q_i,\qquad \Var D=\sum_i q_i(1-q_i).}
\]
又 $N=\sum X_i$，跨类别协方差均为零。由 $I_iX_i=X_i$，
\[
\Cov(I_i,X_i)=\lambda p_i-q_i\lambda p_i
=\lambda p_i e^{-\lambda p_i},
\]
故
\[
\boxed{\Cov(D,N)=\sum_i\lambda p_i e^{-\lambda p_i}.}
\]
\insight 固定总抽样数时各类计数相互制约；随机化为泊松总数后，联合概率恰好分解，覆盖方差也随之简化。

\problem{P32}{由观测值决定的随机和}{首步分解}
$X_i$ 独立服从 $\operatorname{Unif}(0,1)$，给定 $0\le c<1$。在首次 $X_i>c$ 时停止，令 $N$ 为停止次数，$S=\sum_{i=1}^NX_i$。求 $\E S$、$\Var S$ 与 $\Cov(S,N)$。说明为什么不能套用“项数与各项独立”的随机和方差公式。
\solution 记 $q=1-c$，则 $N$ 为参数 $q$ 的几何变量，$\E N=1/q$。由于 $0\le S\le N$，所需二阶矩均有限。设 $(S',N')$ 为独立于第一步的同分布后续过程，则
\[
S=X+\ind\{X\le c\}S',\qquad
N=1+\ind\{X\le c\}N'.
\]
记 $\mu=\E S$、$m_2=\E S^2$。用 $\E X=1/2$、$\E X^2=1/3$ 及 $\E[X\ind\{X\le c\}]=c^2/2$，
\[
\mu=\frac12+c\mu,\qquad
m_2=\frac13+c^2\mu+cm_2.
\]
解得
\[
\boxed{\E S=\frac1{2q},\qquad
\Var S=\frac{6c^2-4c+1}{12q^2}.}
\]
为求混合矩 $h=\E(SN)$，将上述首步等式相乘并取期望：
\[
h=\frac12+\frac{c^2}{2q}+c\mu+ch.
\]
故 $h=(1+c^2)/(2q^2)$，从而
\[
\boxed{\Cov(S,N)=h-\mu/q=\frac{c^2}{2q^2}.}
\]
\pitfall 停止前各项必须不超过 $c$，最后一项必须超过 $c$，因此给定 $N$ 后各项的分布已被改变。独立项数的随机和方差公式不适用；首步分解保留了这种依赖。

\problem{P14}{高维区域体积的阈值}{大数定律与中心极限定理}
记单位立方体内区域的体积为
\[
V_n(\alpha)=\int_{[0,1]^n}
\ind\!\left\{\sum_{i=1}^n x_i^2\le\alpha\sum_{i=1}^nx_i\right\}
\dd x_1\cdots\dd x_n.
\]
对固定实数 $\alpha$，求 $n\to\infty$ 时的极限。进一步令 $\alpha_n=2/3+c/\sqrt n$，其中 $c$ 为固定实数，求 $V_n(\alpha_n)$ 的极限。
\solution 令 $U_i$ 独立均匀分布于 $(0,1)$。体积等于相应事件的概率，而
\[
\E(U_i^2-\alpha U_i)=\frac13-\frac\alpha2.
\]
由大数定律，样本均值趋于此常数。若常数为正，满足不等式的概率趋于零；若为负则趋于一。因此阈值为 $\alpha=2/3$。

阈值处令 $Y_i=U_i^2-2U_i/3$，有 $\E Y_i=0$，且
\[
\sigma^2=\E Y_i^2
=\frac15-\frac13+\frac4{27}=\frac2{135}>0.
\]
中心极限定理给出 $\sum Y_i/\sqrt n\Rightarrow N(0,\sigma^2)$。对变化阈值，原事件等价于
\[
\frac1{\sqrt n}\sum_{i=1}^nY_i
\le c\left(\frac1n\sum_{i=1}^nU_i\right).
\]
右侧由大数定律趋于 $c/2$；极限正态分布连续，因此
\[
\boxed{\lim_{n\to\infty}V_n(\alpha_n)
=\Phi\!\left(c\sqrt{\frac{135}{8}}\right).}
\]
取 $c=0$，固定阈值处极限为 $1/2$。综上，固定 $\alpha$ 的极限依次为 $0$、$1/2$、$1$，分别对应 $\alpha<2/3$、$\alpha=2/3$、$\alpha>2/3$。
\insight 大数定律定位突变阈值；在均值恰好抵消处，必须进一步看 $n^{-1/2}$ 尺度上的波动，不能仅凭极限均值为零判断事件概率。


\par\Needspace{8\baselineskip}
\subsection*{小节练习}
\exercise{OP18}{随机连接面条}
有 $n\ge1$ 根面条，每根有两个端点。每次从尚未连接的端点中等概率选取一对并连接，直到所有 $2n$ 个端点都已两两连接。连接后形成若干闭合环。求环数的期望。

\exercise{OP25}{随机弦的交点数}
在圆周上独立均匀选取 $2n$ 个点，按预先指定的配对构成 $n\ge1$ 条弦。求圆内部不同交点个数的期望。

\exercise{OP26}{抽齐四张A}
将一副不含大小王的 $52$ 张扑克牌均匀洗牌，逐张无放回抽取，直到抽到全部 $4$ 张 A 为止。设抽取张数为 $T$，包含最后抽到的那张 A。求 $\E T$。


\section{等待时间}
\noindent\textbf{要点与方法.}\enspace
等待问题先明确停止是否包含触发事件的那次观测，以及是在每步还是每轮末检查。对非负整数值的停止次数 $T$，尾概率直接给出
\[
\E T=\sum_{n\ge0}\Prob(T>n),
\]
等式也允许两边为无穷；完整分布则由相邻尾概率相减得到。

独立均匀比例的负对数是指数增量，可把乘法缩减转成到达时间问题。仅涉及相邻大小关系时，独立同分布连续样本的等可能排名往往比密度更方便。正增量的越界问题可汇总各次累计和的密度。

写出一、二阶矩递推前，应确认所用矩有限；处理无限随机和或乘积的极限时，则要说明交换极限与期望的依据。
\par\smallskip



\problem{P06}{未知阈值下的停止次数}{混合分布与后验}
独立变量 $A,U_1,U_2,\ldots$ 均服从 $\operatorname{Unif}(0,1)$，其中阈值 $A$ 不可见。在累计和首次超过 $A$ 时停止，次数记为 $N$。求 $N$ 的分布、均值和方差；只观察到 $N=n$ 时，求阈值 $A$ 的后验密度与均值。
\solution 给定 $A=a$，前 $k$ 个增量之和不超过 $a\le1$ 的事件是边长约束不再起作用的单纯形，其体积为 $a^k/k!$。因此
\[
\Prob(N>k)=\int_0^1\frac{a^k}{k!}\dd a=\frac1{(k+1)!},\quad k\ge0.
\]
相邻尾概率相减，得到
\[
\boxed{\Prob(N=n)=\frac n{(n+1)!},\quad n\ge1.}
\]
由整数变量的尾和恒等式，
\[
\E N=\sum_{k\ge0}\frac1{(k+1)!}=e-1,
\quad \E N^2=\sum_{k\ge0}\frac{2k+1}{(k+1)!}=e+1.
\]
故 $\boxed{\Var N=3e-e^2}$。

对后验，给定阈值的似然为
\[
\Prob(N=n\mid A=a)=\frac{a^{n-1}}{(n-1)!}-\frac{a^n}{n!}.
\]
乘以均匀先验密度并除以 $\Prob(N=n)$，得
\[
\boxed{f_{A\mid N=n}(a)=(n+1)a^{n-1}-\frac{n+1}{n}a^n,
\quad0<a<1.}
\]
再积分求均值：
\[
\boxed{\E(A\mid N=n)=1-\frac{n+1}{n(n+2)}.}
\]
例如一次就停止时，阈值的后验均值为 $1/3$；停止次数越多，观察越支持较大的阈值。

\problem{P20}{每轮两次缩小的越界时间}{指数等待与奇偶分解}
每轮把当前长度乘以两个新的独立 $\operatorname{Unif}(0,1)$ 变量，所有变量相互独立。从 $R_0=1$ 出发，令 $R_n=\prod_{i=1}^nU_iV_i$，在首次 $R_n<x$ 时停止，$0<x<1$，次数为 $N$。求 $N$ 的分布、均值，以及终点比例 $Y=R_N/x$ 的密度。
\solution 记 $t=-\log x$。所有负对数增量都是参数为 $1$ 的独立指数变量。把每轮的两个增量依次展开，令 $K$ 为到时刻 $t$ 已完成的指数等待次数。对 $k\ge1$，以第 $k$ 次到达时刻 $s$ 积分，
\[
\Prob(K=k)=\int_0^t\frac{s^{k-1}e^{-s}}{(k-1)!}e^{-(t-s)}\dd s
=e^{-t}\frac{t^k}{k!};
\]
$k=0$ 时概率为 $e^{-t}$，故 $K\sim\operatorname{Pois}(t)$。

只有每两个增量结束后才检查长度，所以 $N=\lfloor K/2\rfloor+1$，即
\[
\boxed{\Prob(N=n)=e^{-t}\left[\frac{t^{2n-2}}{(2n-2)!}
+\frac{t^{2n-1}}{(2n-1)!}\right],\quad n\ge1.}
\]
由泊松生成函数，$\Prob(K\text{ 为奇数})=(1-e^{-2t})/2$。于是
\[
\boxed{\E N=1+\frac{\E K-\Prob(K\text{ 为奇数})}{2}
=\frac{3+2t+e^{-2t}}4.}
\]
由指数分布的无记忆性，从 $t$ 到下一次检查的距离 $O=-\log Y$，在 $K$ 为奇数时是一个指数等待，在 $K$ 为偶数时是两个指数等待之和。记 $w_o=(1-e^{-2t})/2$、$w_e=(1+e^{-2t})/2$，则
\[
f_O(o)=e^{-o}(w_o+w_eo),\quad o>0.
\]
作 $y=e^{-o}$ 的变量变换，得到
\[
\boxed{f_Y(y)=w_o-w_e\log y,\quad0<y<1.}
\]
\insight 批量检查改变了越界后的剩余等待。是否已经走完半轮决定还需一个还是两个指数增量，因此终点比例不再均匀。

\problem{P21}{两种整数步长的更新过程}{递推与越界量}
独立增量 $X_i$ 取值 $1,2$，概率分别为 $1-q,q$，其中 $0\le q\le1$。对整数 $m\ge1$，令 $T_m=\min\{n:\sum_{i=1}^nX_i\ge m\}$，越界量为 $O_m=\sum_{i=1}^{T_m}X_i-m$。求 $\E T_m$ 和 $O_m$ 的分布。
\solution 记 $e_m=\E T_m$，并置 $e_0=e_{-1}=0$。按第一步条件化，
\[
e_m=1+(1-q)e_{m-1}+qe_{m-2},\quad m\ge1.
\]
令 $d_m=e_m-e_{m-1}$，则 $d_1=1$，且 $d_m=1-qd_{m-1}$。由等比数列求和，
\[
d_m=\frac{1-(-q)^m}{1+q},\qquad
\boxed{e_m=\frac m{1+q}+\frac{q[1-(-q)^m]}{(1+q)^2}.}
\]
每步至少前进 $1$，故 $T_m\le m$。写成有限指示变量和，利用 $\{T_m\ge i\}$ 只依赖前 $i-1$ 个增量，得到
\[
\E\sum_{i=1}^{T_m}X_i
=\sum_{i=1}^m\E[X_i\ind\{T_m\ge i\}]
=(1+q)e_m.
\]
越界量仅能取 $0,1$，所以
\[
\boxed{\Prob(O_m=1)=\E O_m
=\frac{q[1-(-q)^m]}{1+q},\quad
\Prob(O_m=0)=1-\Prob(O_m=1).}
\]
$q=0$ 时恰需 $m$ 步；$q=1$ 时恰需 $\lceil m/2\rceil$ 步，越界量由 $m$ 的奇偶性决定，均与公式相符。
\insight 整数步长留下了交替项 $(-q)^m$。这项反映格点结构，不能用仅含平均步长的近似替代。

\problem{P22}{第一次出现局部峰值}{随机排名与尾概率}
$X_1,X_2,\ldots$ 独立同分布于任意连续分布。令
\[
N=\min\{n\ge3:X_{n-2}<X_{n-1}>X_n\},
\]
即在确认第一个局部峰值时停止，包含峰值后的下降观测。求 $N$ 的分布、均值和方差。
\solution 对前 $n$ 个互异观测，没有局部峰值等价于相邻比较符号中不出现“先升后降”。因此整个序列只能先递减、再递增，中间允许某一段为空。

在 $n$ 个排名中，最小者作为谷底。其余 $n-1$ 个排名可各自选择放在谷底左边或右边；左边必须递减、右边必须递增，选定两侧后次序唯一。因此无峰值的排列恰有 $2^{n-1}$ 个。连续独立同分布保证 $n!$ 个排名等可能，故
\[
\boxed{\Prob(N>n)=\frac{2^{n-1}}{n!},\quad n\ge1,}
\]
另有 $\Prob(N>0)=1$。对 $n\ge3$，
\[
\Prob(N=n)=\frac{2^{n-2}}{(n-1)!}-\frac{2^{n-1}}{n!}.
\]
尾概率趋于零，说明停止几乎必然发生。由尾和公式及指数级数，
\[
\E N=1+\frac12\sum_{n\ge1}\frac{2^n}{n!}
=\frac{e^2+1}{2},
\]
\[
\E N^2=1+\frac12\sum_{n\ge1}(2n+1)\frac{2^n}{n!}
=\frac{5e^2+1}{2}.
\]
因此
\[
\boxed{\Var N=\frac{5e^2+1}{2}-\left(\frac{e^2+1}{2}\right)^2.}
\]
\insight 不必知道观测值的密度，只需数满足相邻不等式的排名。停止时刻是否包含确认峰值的最后一次观测，必须在题设中明确。

\problem{P23}{越界前后的两个间隙}{更新密度}
$U_i$ 独立服从 $\operatorname{Unif}(0,1)$，记 $S_k=\sum_{i=1}^kU_i$，$S_0=0$，$N=\min\{k:S_k>1\}$。令 $A=1-S_{N-1}$ 为越界前距离，$O=S_N-1$ 为越界量。求 $(A,O)$ 的联合密度、最后一步长度 $J=U_N$ 的密度，以及 $\Prob(O>A)$。
\solution 对 $0<s<1$，第 $k$ 个正累计和的密度为 $s^{k-1}/(k-1)!$。所有正累计和的密度相加得到
\[
u(s)=\sum_{k\ge1}\frac{s^{k-1}}{(k-1)!}=e^s.
\]
第一步不可能超过 $1$，故最后越界前的累计和必为正，没有来自 $S_0$ 的原子。若越界前在 $s$，下一步为 $v$，则联合密度贡献为 $e^s\dd s\dd v$，条件是 $s+v>1$。由于增量为正，这个跨越步骤唯一。

变换 $a=1-s$、$o=s+v-1$ 的雅可比绝对值为 $1$，所以
\[
\boxed{f_{A,O}(a,o)=e^{1-a},\quad 0<a<1,\quad0<o<1-a.}
\]
积分 $\int_0^1e^{1-a}(1-a)\dd a=1$，确为概率密度。最后一步 $J=A+O$，因此
\[
\boxed{f_J(j)=\int_0^je^{1-a}\dd a=e-e^{1-j},\quad0<j<1.}
\]
要有 $O>A$，必须 $a<1/2$，故
\[
\boxed{\Prob(O>A)=\int_0^{1/2}e^{1-a}(1-2a)\dd a
=2\sqrt e-e.}
\]
\insight 把所有可能的越界前累计和相加，再接上最后一个增量，就能得到越界前后两个量的联合分布，而无须分别枚举停止次数。

\problem{P24}{递增过程的总剩余距离}{分布递推}
给定 $R_0=x\in[0,1)$，给定过去的观测，每次使用新的独立随机数，在区间 $(R_{j-1},1)$ 上条件均匀选取 $R_j$。令
\[
Q=\sum_{j\ge1}(1-R_j).
\]
证明 $Q$ 几乎处处有限，求其均值与方差，并求 $D=Q/(1-x)$ 的拉普拉斯变换 $L(t)=\E e^{-tD}$，$t\ge0$。
\solution 令 $U_j$ 为独立均匀变量，则剩余距离可写为
\[
1-R_j=(1-x)\prod_{i=1}^jU_i.
\]
非负项允许交换求和与期望，故 $\E Q=(1-x)\sum_{j\ge1}2^{-j}=1-x$。非负随机变量具有有限期望，因而 $Q<\infty$ 几乎处处。

归一化后 $D=U(1+D')$，其中 $D'$ 与 $D$ 同分布且独立于 $U$。二阶矩确实有限：由三角不等式，各有限部分和的 $L^2$ 范数不超过 $\sum_{j\ge1}3^{-j/2}<\infty$，再用单调收敛。于是
\[
\E D=1,\qquad
\E D^2=\frac13(1+2\E D+\E D^2)=\frac32.
\]
因此
\[
\boxed{\E Q=1-x,\qquad \Var Q=\frac{(1-x)^2}{2}.}
\]
对 $t>0$，分布递推给出
\[
L(t)=\int_0^1e^{-tu}L(tu)\dd u,
\qquad tL(t)=\int_0^te^{-s}L(s)\dd s.
\]
右侧可微，且 $L(t)>0$，所以
\[
\frac{L'(t)}{L(t)}=\frac{e^{-t}-1}{t}.
\]
用 $L(0)=1$ 积分，得到
\[
\boxed{L(t)=\exp\!\left\{\int_0^t\frac{e^{-s}-1}{s}\dd s\right\}.}
\]
被积函数在零点连续延拓为 $-1$，因此该表达式在 $t=0$ 也有定义。
\insight 剩余距离具有乘法结构，但所有剩余距离之和满足简单的加法递推。先证明矩有限，再对分布递推取矩，避免用未证实的有限量作代数运算。


\par\Needspace{8\baselineskip}
\subsection*{小节练习}
\exercise{OP22}{交替升降的停止时刻}
依次独立生成 $U_1,U_2,\ldots\sim\operatorname{Unif}(0,1)$，要求
\[
U_1<U_2>U_3<U_4>U_5<\cdots.
\]
首次生成一个使上述交替不等式失效的数时停止。记总生成次数为 $T$，包括导致失效的最后一次。求 $\E T$。

\exercise{OP24}{递增均匀序列的乘积}
先取 $R_1\sim\operatorname{Unif}(0,1)$；之后给定过去，$R_{k+1}$ 在 $(R_k,1)$ 上均匀分布，每一步使用新的独立随机数。求
\[
\lim_{n\to\infty}\E\!\left[\prod_{k=1}^nR_k\right].
\]


\section{分布与几何}
\noindent\textbf{要点与方法.}\enspace
求比值、极值和相对位置的分布，可从联合密度作变量变换，写清取值域与雅可比。条件缩放有时能消去已知位置，直接证明独立性。均匀切点产生的间距在单纯形上具有常密度，但各间距并不独立。

几何概率必须指定位置与方向的抽样方式。距离阈值可转为邻域内的二项计数，交点与配对可用指示变量求均值；共点或边界情形须单独辨认。分布对称也不保证期望存在。

随机着色可用非负整数坏事件数的期望小于一证明存在性。保密问题则要比较给定消息的密文分布，不能仅凭密文边际均匀判断联合保密。
\par\smallskip



\problem{P04}{给定乘积后的条件分布}{变量变换与独立性}
$U,V$ 独立服从 $\operatorname{Unif}(0,1)$，令 $P=UV$、$Q=\log U/\log P$。证明 $P,Q$ 独立，并求给定 $P=p\in(0,1)$ 时 $U$ 的条件密度及条件均值。比较该均值与 $\sqrt p$。
\solution 令 $X=-\log U$、$Y=-\log V$，则二者是独立的参数为 $1$ 的指数变量。对 $S=X+Y$、$Q=X/S$，逆变换为 $X=QS$、$Y=(1-Q)S$，雅可比为 $s$。联合密度为
\[
f_{S,Q}(s,q)=s e^{-s},\quad s>0,\quad0<q<1.
\]
它分解为 $S$ 的伽马密度与 $Q$ 的均匀密度，故 $Q$ 与 $S$ 独立，进而与 $P=e^{-S}$ 独立。还得到 $f_P(p)=-\log p$。

给定 $P=p$，有 $U=p^Q$，其中 $Q$ 仍均匀。因 $q\mapsto p^q$ 递减，对 $p<u<1$，
\[
\Prob(U\le u\mid P=p)=1-\frac{\log u}{\log p}
=\frac{\log(u/p)}{\log(1/p)}.
\]
因此
\[
\boxed{f_{U\mid P=p}(u)=\frac1{u\log(1/p)},\quad p<u<1,}
\]
且
\[
\boxed{\E(U\mid P=p)=\int_0^1p^q\dd q
=\frac{1-p}{\log(1/p)}>\sqrt p.}
\]
最后的不等式来自严格凸函数 $q\mapsto p^q$ 的詹森不等式，因为 $\E Q=1/2$ 且 $Q$ 非退化。
\pitfall 给定连续变量的精确取值需要条件密度，不能把零概率事件直接代入离散贝叶斯公式。乘积固定也不意味着两个因子各等于其几何平均。

\problem{P08}{两个高斯量的相对差}{旋转与极坐标}
$X,Y$ 为独立标准正态变量，令
\[
R=\frac{X-Y}{X+Y},\qquad H=X^2+Y^2.
\]
分母为零的事件概率为零。求 $R$ 的密度，证明 $R,H$ 独立，并对 $c\ge0$ 求 $\Prob(|R|\le c)$。$R$ 的均值是否等于零？
\solution 令 $D=(X-Y)/\sqrt2$、$S=(X+Y)/\sqrt2$。正交变换保留联合标准正态密度，所以 $D,S$ 仍独立标准正态，且 $H=D^2+S^2$。写
\[
S=\rho\cos\theta,\qquad D=\rho\sin\theta.
\]
极坐标联合密度为 $(2\pi)^{-1}\rho e^{-\rho^2/2}$，可分解为半径密度与均匀角度密度。因此角度与半径独立，$R=\tan\theta$ 与 $H=\rho^2$ 独立。

将角度按模 $\pi$ 化到 $(-\pi/2,\pi/2)$，其密度为 $1/\pi$。用 $\theta=\arctan r$ 变换得
\[
\boxed{f_R(r)=\frac1{\pi(1+r^2)},\quad r\in\mathbb R,}
\qquad
\boxed{\Prob(|R|\le c)=\frac2\pi\arctan c.}
\]
尽管密度关于零对称，$\int_0^\infty r/[\pi(1+r^2)]\dd r=\infty$，正负部分的期望均发散。因此 $\E R$ 不存在，不能写成零。
\insight 旋转把相对差化为两个独立高斯量的比值；极坐标再把方向与大小分开。分布对称与均值存在是两个不同条件。

\problem{P10}{相关二进制选择的极限}{马尔可夫相关结构}
从单位区间开始，第 $j$ 步按 $B_j=0$ 或 $1$ 选择当前区间的左半或右半。$B_1$ 等概率取 $0,1$；以后以概率 $p\in[0,1]$ 保持前一位，以概率 $1-p$ 翻转，条件于当前位后的转移与过去无关。令极限位置为 $X$。求 $\E X$、$\Var X$，并说明 $p=0,1/2,1$ 的分布及前 $n$ 位近似的确定性误差界。
\solution 区间长度每步减半，唯一极限为
\[
X=\sum_{j\ge1}2^{-j}B_j.
\]
链处于平稳分布，故 $\E B_j=1/2$。记 $Z_j=B_j-1/2$、$\lambda=2p-1$，一步条件期望为 $\E(Z_{j+1}\mid Z_j)=\lambda Z_j$。迭代得到
\[
\Cov(B_i,B_j)=\frac14\lambda^{|i-j|}.
\]
绝对收敛允许交换求和；因为 $\sum_i4^{-i}=1/3$，
\[
\begin{aligned}
\Var X
&=\frac14\sum_i4^{-i}
+\frac12\sum_{k\ge1}\lambda^k\sum_{i\ge1}2^{-i}2^{-(i+k)}\\
&=\frac1{12}+\frac16\sum_{k\ge1}(\lambda/2)^k
=\boxed{\frac{2+\lambda}{12(2-\lambda)}}.
\end{aligned}
\]
同时 $\boxed{\E X=1/2}$。

$p=0$ 时数字交替，$X$ 等概率为 $1/3,2/3$；$p=1$ 时数字全部相同，$X$ 等概率为 $0,1$；$p=1/2$ 时数字独立，每个长度为 $2^{-n}$ 的二进区间概率都是 $2^{-n}$，从而 $X$ 均匀分布。$p=0,1/2,1$ 时，方差依次为 $1/36,1/12,1/4$。

令 $X_n=\sum_{j=1}^n2^{-j}B_j$，则总有 $0\le X-X_n\le2^{-n}$；若用所选区间中点近似，绝对误差不超过 $2^{-n-1}$。
\insight 边际上每位都公平，并不足以推出极限位置均匀；位与位之间的相关性会通过几何加权累积到方差中。

\problem{P11}{拉普拉斯样本的中位数误差}{尾积分与估计比较}
$X_1,X_2,X_3$ 独立服从位置为 $c$、尺度为 $b>0$ 的拉普拉斯分布，密度为 $e^{-|x-c|/b}/(2b)$。用样本中位数 $M$ 或样本均值 $\overline X$ 估计 $c$。证明二者无偏，分别求均方误差，并判断哪一个更小。
\solution 联合分布在关于 $c$ 的反射下不变，故 $M-c$ 的分布对称；又 $|M-c|\le\sum_i|X_i-c|$，其期望有限，因而 $\E M=c$。单个样本由对称性亦有均值 $c$，故 $\E\overline X=c$。

对 $t\ge0$，记 $q(t)=\Prob(X_i>c+t)=e^{-t/b}/2$。中位数大于 $c+t$ 要求至少两个样本超过该值，故
\[
\Prob(M>c+t)=3q(t)^2-2q(t)^3.
\]
利用左右对称性，
\[
\Prob(|M-c|>t)=\frac32e^{-2t/b}-\frac12e^{-3t/b}.
\]
平方的尾积分给出
\[
\begin{aligned}
\E(M-c)^2
&=2\int_0^\infty t\left(\frac32e^{-2t/b}-\frac12e^{-3t/b}\right)\dd t\\
&=\frac{3b^2}{4}-\frac{b^2}{9}
=\boxed{\frac{23b^2}{36}}.
\end{aligned}
\]
而 $\Var X_i=b^{-1}\int_0^\infty t^2e^{-t/b}\dd t=2b^2$，独立性给出
\[
\boxed{\E(\overline X-c)^2=\frac{2b^2}{3}=\frac{24b^2}{36}.}
\]
因此在这个样本量和分布下，中位数的均方误差小 $b^2/36$。
\insight 比较估计量须指定分布和损失。中位数的尾概率可以直接用二项计数写出，再用尾积分计算平方误差，避免繁琐的分段密度积分。

\problem{P35}{剩余区间中的次序位置}{条件缩放}
从 $(0,1)$ 独立均匀抽取 $n$ 个样本，次序统计量记为 $U_{(1)}<\cdots<U_{(n)}$。对 $1\le i<j\le n$，令
\[
W=\frac{U_{(j)}-U_{(i)}}{1-U_{(i)}}.
\]
证明 $W$ 与 $U_{(i)}$ 独立，求 $W$ 的分布、均值和方差，并计算 $\Prob(U_{(j)}>(1+U_{(i)})/2)$。
\solution 给定 $U_{(i)}=u$，上方 $n-i$ 个样本的无序集合服从 $(u,1)$ 上独立均匀样本的规律。这可由次序统计量的联合密度为常数，并在 $u$ 两侧分别分解积分区域得到。将上方样本缩放为 $(v-u)/(1-u)$ 后，便是 $m=n-i$ 个标准均匀样本。

$W$ 为其中第 $r=j-i$ 个次序统计量。通过在 $w$ 下方放 $r-1$ 个样本、在 $w$ 附近放一个样本、其余放在上方，得到
\[
f_W(w)=\frac{m!}{(r-1)!(m-r)!}
 w^{r-1}(1-w)^{m-r},\quad0<w<1.
\]
这个条件密度不依赖 $u$，所以 $W$ 与 $U_{(i)}$ 独立，且 $W\sim\operatorname{Beta}(r,m-r+1)$。用贝塔积分计算一、二阶矩，得到
\[
\boxed{\E W=\frac r{m+1},\qquad
\Var W=\frac{r(m-r+1)}{(m+1)^2(m+2)}.}
\]
$W>1/2$ 等价于 $m$ 个缩放样本中至多 $r-1$ 个落在左半区间，因此
\[
\boxed{\Prob\!\left(U_{(j)}>\frac{1+U_{(i)}}2\right)
=2^{-m}\sum_{k=0}^{r-1}\binom mk.}
\]
\insight 绝对位置互相依赖，条件缩放后的相对位置却可以独立。找出条件密度中能被消去的尺度，比直接算协方差更能揭示结构。

\problem{P36}{最短随机段与超长段计数}{单纯形平移}
用 $n-1$ 个独立均匀切点把单位线段分成 $n\ge2$ 段，长度为 $D_1,\ldots,D_n$。令 $Z=\min_iD_i$，并对 $t\ge0$ 令 $K_t=\sum_i\ind\{D_i>t\}$。求 $Z$ 的分布、均值和方差，以及 $K_t$ 的概率生成函数与均值。
\solution 切点排序后的联合密度恒为 $(n-1)!$。由切点到间距的线性变换可知，间距向量均匀分布于单纯形 $D_i\ge0,\sum D_i=1$。

若指定 $r$ 个间距均超过 $t$，从这 $r$ 个坐标分别减去 $t$，剩余单纯形的总长度变为 $1-rt$。因其维数为 $n-1$，该事件概率为 $(1-rt)_+^{n-1}$。取 $r=n$，
\[
\boxed{\Prob(Z>z)=(1-nz)^{n-1},\quad0\le z\le1/n.}
\]
支持集之外由 $Z\in[0,1/n]$ 确定。用尾积分及代换 $u=nz$，
\[
\E Z=\frac1{n^2},\qquad
\E Z^2=\frac2{n^3(n+1)},\qquad
\boxed{\Var Z=\frac{n-1}{n^4(n+1)}.}
\]
对计数，展开指示变量乘积：
\[
z^{K_t}=\prod_{i=1}^n[1+(z-1)\ind\{D_i>t\}].
\]
每个 $r$ 元子集有相同的联合概率，故
\[
\boxed{\E z^{K_t}=\sum_{r=0}^n\binom nr(z-1)^r(1-rt)_+^{n-1}.}
\]
此处 $r=0$ 项为 $1$。对 $z$ 在 $1$ 处求导，得到
\[
\boxed{\E K_t=n(1-t)_+^{n-1}.}
\]
\insight 对若干间距同时设下界，等价于平移并缩小同一个单纯形。一次联合概率计算便能同时处理最短段和超阈值段数。

\problem{P37}{固定位置的近邻距离}{邻域计数}
在单位线段上独立均匀抽取 $n\ge1$ 个样本点，另固定一个观测位置 $a\in[0,1]$，它不属于随机样本。把样本到 $a$ 的距离从小到大排列，记第 $k$ 个为 $D_{(k)}$，$1\le k\le n$。求一般 $a$ 下 $D_{(k)}$ 的分布，并在 $a=0$ 与 $a=1/2$ 时给出密度、均值和方差。
\solution 对 $r\ge0$，观测位置半径 $r$ 的邻域长度为
\[
\ell_a(r)=\min(1,a+r)-\max(0,a-r).
\]
这里当 $r$ 超过所有可能距离时，长度为 $1$。每个样本独立以概率 $\ell_a(r)$ 落入邻域，故邻域样本数为参数 $n,\ell_a(r)$ 的二项变量。第 $k$ 近距离超过 $r$，恰好等价于邻域中至多 $k-1$ 个样本。因此
\[
\boxed{\Prob(D_{(k)}>r)=\sum_{j=0}^{k-1}\binom nj
\ell_a(r)^j[1-\ell_a(r)]^{n-j},\quad r\ge0.}
\]
当 $a=0$ 时，单个距离本身均匀分布于 $(0,1)$，于是
\[
f_{D_{(k)}}(r)=\frac{n!}{(k-1)!(n-k)!}r^{k-1}(1-r)^{n-k},\quad0<r<1.
\]
这是 $\operatorname{Beta}(k,n-k+1)$ 密度，贝塔积分给出
\[
\boxed{\E D_{(k)}=\frac{k}{n+1},\qquad
\Var D_{(k)}=\frac{k(n-k+1)}{(n+1)^2(n+2)}.}
\]
当 $a=1/2$ 时，单个距离的分布函数在 $[0,1/2]$ 上为 $2r$，故 $2D_{(k)}$ 服从同一贝塔分布。相应密度为
\[
f_{D_{(k)}}(r)=\frac{2n!}{(k-1)!(n-k)!}(2r)^{k-1}(1-2r)^{n-k},\quad0<r<1/2,
\]
均值为端点情形的一半，方差为四分之一。
\insight 把第 $k$ 近距离超过阈值的事件转成固定邻域中的计数事件，可统一处理任意位置和任意次序。端点与中点的区别来自邻域长度的增长速度，而不是样本之间的互相配对。

\problem{P38}{中间排名的零和奖金}{条件排名与收益}
三名参与者从 $1,\ldots,M$ 中不放回均匀各取一个数，$M\ge3$。排序后为 $A<B<C$。中间者获得奖金 $C-A$，两端者各支付 $(C-A)/2$。已知你抽到 $x$，求你的条件期望净收益 $g(x)$，并确定最有利的 $x$。计算 $M=9$ 时的最优数值和收益。
\solution 另两人的无序抽数对在 $D=\binom{M-1}{2}$ 个可能中均匀分布。若你居中，一低一高的配对数为 $(x-1)(M-x)$；两侧条件均匀，故平均极差为
\[
\frac{x+1+M}{2}-\frac x2=\frac{M+1}{2}.
\]
若你最小，另两数都在 $x+1,\ldots,M$，共有 $\binom{M-x}{2}$ 对。设 $L=M-x$；从 $1,\ldots,L$ 取两个数，其最大值均值为 $2(L+1)/3$，因为
\[
\sum_{b=2}^Lb(b-1)\big/\binom L2=\frac{2(L+1)}3.
\]
因此你的条件净收益为 $-(M-x+1)/3$。对你最大的情形，对称地得到条件净收益 $-x/3$，配对数为 $\binom{x-1}{2}$。不可能的排名对应配对数为零，不需赋予其条件均值。

合并三类，
\[
\begin{aligned}
g(x)=\frac1D\bigg[&\frac{M+1}{2}(x-1)(M-x)\\
&-\frac{M-x+1}{3}\binom{M-x}{2}
-\frac{x}{3}\binom{x-1}{2}\bigg].
\end{aligned}
\]
化简为
\[
\boxed{g(x)=\frac{2M(x-1)(M-x)}{(M-1)(M-2)}-\frac M3.}
\]
乘积 $(x-1)(M-x)$ 在最接近 $(M+1)/2$ 的整数处最大。$M$ 为奇数时最优值唯一，为偶数时中间两个整数并列。$M=9$ 时最优 $x=5$，$g(5)=15/7$。
\insight 条件收益同时取决于排名概率和该排名下的平均极差。这里每组净收益之和为零，但知道自己抽到中间数后仍有正的条件期望收益。


\problem{CP01}{跨越平行线的随机针}{几何概率与交点数}
平面上画间距为 $h>0$ 的平行直线。投下一根长度为 $l>0$ 的针，假设针中心到最近直线的距离在 $[0,h/2]$ 上均匀，针与平行线的锐角在 $[0,\pi/2]$ 上均匀，且二者独立。求针至少碰到一条线的概率和针与直线的交点数期望；针长可以超过线距。
\solution 记角度为 $\Theta$、中心距离为 $D$、长度比为 $r=l/h$。针的法向投影长度为 $l\sin\Theta$，碰线等价于 $D\le l\sin\Theta/2$。条件化于角度，
\[
p(r)=\frac2\pi\int_0^{\pi/2}\min(1,r\sin\theta)\dd\theta.
\]
若 $r\le1$，积分直接给出 $p(r)=2r/\pi$。若 $r>1$，以 $\theta_0=\arcsin(1/r)$ 分开积分，得到
\[
\boxed{p(r)=\begin{cases}
2r/\pi,&0<r\le1,\\
\dfrac2\pi\left(r-\sqrt{r^2-1}+\arccos(1/r)\right),&r>1.
\end{cases}}
\]
再求交点数 $J$。给定角度，针投影为长度 $b=l\sin\theta$ 的区间，直线位置构成间距 $h$ 的格点。交点数在中心位置关于最近直线反射时不变，所以距离均匀的模型可等价地用中心在一个周期内均匀平移来计算。此时，任意长度 $b$ 的区间所覆盖格点数的均值为 $b/h$：把各格点的被覆盖指示函数积分并相加，恰好覆盖长度为 $b$ 的平移范围。故
\[
\boxed{\E J=\E\left(\frac{l\sin\Theta}{h}\right)=\frac{2l}{\pi h}.}
\]
当 $l\le h$ 时，几乎必然至多一个交点，故概率与期望相等；长针可能跨多条线，两者必须区分。
\insight “随机投针”必须指定位置和方向的分布。先条件化角度，再求几何长度比例，可以同时处理短针和长针，而不依赖图形直觉。

\problem{CP02}{避免单色集合的二色着色}{概率方法与条件期望}
有限点集 $V$ 上给定 $m$ 个子集，每个子集至少含 $k\ge2$ 个点。若 $m<2^{k-1}$，证明存在一种红蓝二色着色，使每个指定子集都同时含两种颜色。进一步说明如何逐点确定颜色，而不必穷举全部着色。
\solution 先独立等概率把各点染成红或蓝。含 $s$ 个点的子集全红或全蓝的概率为 $2^{1-s}\le2^{1-k}$。令 $Z$ 为单色指定子集的个数，则
\[
\E Z\le m2^{1-k}<1.
\]
若每种着色都有单色子集，必有 $Z\ge1$，与期望小于 $1$ 矛盾。因此存在 $Z=0$ 的着色。联合界还给出成功概率至少为 $1-m2^{1-k}>0$。

可把这一存在性证明变成逐点选择。每次已有部分颜色固定时，计算当前条件期望。下一个点随机取红、蓝的两种条件期望的平均，等于当前值，因此至少一种颜色不会使条件期望增加，固定为该颜色即可。

这些条件期望可以直接计算。一个子集若已出现两色，单色概率为零；若已固定点均同色、尚有 $u$ 个未染点，单色概率为 $2^{-u}$；若尚无点固定，概率仍为 $2^{1-s}$。将各子集概率相加即可。所有点固定后，条件期望就是整数 $Z$，仍小于 $1$，故必为零。
\insight 随机性不仅用于计算概率，也能证明确定对象存在。每一步保留较小的条件期望，就把随机证明转成了确定选择规则；这里给出的数量条件是充分条件，不必是必要条件。

\problem{CP03}{一次一密与密钥复用}{条件独立与信息泄漏}
消息与密钥取值于 $\mathbb Z_q=\{0,\ldots,q-1\}$，$q\ge2$，加密为 $C=M+K\pmod q$。密钥 $K$ 与消息 $M$ 独立。证明均匀密钥使 $C$ 与 $M$ 独立；若消息每个值均有正概率，证明完全保密也要求密钥均匀。再令两条消息独立均匀，复用同一个独立均匀密钥得到 $C_1,C_2$，求给定两份密文后消息对的条件分布，并指出泄漏的确定信息。
\solution 若密钥均匀，对任意 $m,c$，
\[
\Prob(C=c\mid M=m)=\Prob(K=c-m)=\frac1q.
\]
条件分布不依赖消息，所以密文均匀且独立于消息，消息先验不必均匀。

反过来，若密文独立于消息且每个消息值均有正概率，则对固定 $c$，$\Prob(K=c-m)=\Prob(C=c)$ 对所有 $m$ 相等。随着 $m$ 遍历 $\mathbb Z_q$，$c-m$ 遍历所有密钥值，因而各密钥概率相等，只能都为 $1/q$。满支持假设不可省略，例如恒定消息与任何密文都独立。

复用时，对给定密文 $(c_1,c_2)$，每个密钥 $k$ 唯一对应消息对 $(c_1-k,c_2-k)$。原先每个三元组 $(M_1,M_2,K)$ 概率为 $q^{-3}$，所以该密文概率为 $q\cdot q^{-3}=q^{-2}$。于是
\[
\boxed{\begin{aligned}
&\Prob(M_1=m_1,M_2=m_2\mid C_1=c_1,C_2=c_2)\\
&\qquad=\begin{cases}
1/q,&m_2-m_1=c_2-c_1\pmod q,\\
0,&\text{其他情形}.
\end{cases}\end{aligned}}
\]
单条消息的后验仍均匀，但差值 $M_2-M_1$ 已被密文差值完全暴露。消息对从 $q^2$ 个可能缩减为 $q$ 个，不能再称为联合完全保密。
\insight 单份密文的保密性不保证多份密文的联合保密性。重复使用密钥后，相减消去密钥，正是依赖结构泄漏信息的原因。


\par\Needspace{8\baselineskip}
\subsection*{小节练习}
\exercise{OP35}{样本最大值与最小值的相关性}
设 $X_1,\ldots,X_n$ 独立同分布于 $\operatorname{Unif}(0,1)$，其中 $n\ge2$。令 $L=\min_iX_i$、$H=\max_iX_i$，求相关系数 $\Corr(L,H)$。

\exercise{OP36}{随机断棒组成多边形}
在单位线段上独立均匀选取 $n-1$ 个切点，其中 $n\ge3$，由此得到 $n$ 段。求这 $n$ 段能够作为一个非退化 $n$ 边形各边的概率。切点均在原单位线段上选取，不是每次先随机选一段再切。

\exercise{OP37}{互为最近邻的点数}
在 $(0,1)$ 中独立均匀抽取 $N\ge3$ 个点。每个点的邻居是其他点中距离它最近的点；距离相同的零概率情形可任意处理。若某点的邻居也把它选作邻居，称该点为好点。求好点数量的期望。


\section{练习题}
\noindent\textbf{要点与方法.}\enspace
这些练习把前面的工具用于新的观测方式和取样规则。条件概率题应写出完整记录的似然，并检验缺失或筛选是否改变分布；条件期望给出平方损失下的最优预测，但需有有限二阶矩。

计数方差可按两组对象的交集分类；停止成本可考察“本次抽样会发生”是否只依赖过去，再利用新增量与过去独立。分布题要区分等概率选段与按长度选段，并核对端点、归一化和尺度变化。联合符号事件也不能仅凭边际概率判断独立。
\par\smallskip



\exercise{PG1-1}{只记录是否出现正面}
以概率 $\pi$ 选择正面概率为 $p_1$ 的硬币，否则选择正面概率为 $p_2$ 的硬币，其中 $0<\pi,p_1,p_2<1$。选定后独立抛掷同一枚硬币 $n\ge1$ 次。记录只告诉你“至少出现过一次正面”，没有正面个数或顺序。求硬币类型的后验概率和下一次正面的预测概率。说明为什么不能把这一记录当作一次正面观测。

\exercise{PG1-2}{期限内的竞争时钟}
独立时钟 $A,B$ 的响铃时间分别服从参数为 $\alpha,\beta>0$ 的指数分布。记首次响铃时间为 $T$。已知 $A$ 先响，并且 $T<t$，其中 $t>0$ 固定。求 $T$ 的条件密度和均值，并讨论 $t\to\infty$ 的极限。

\exercise{PG1-3}{含缺失报告的后验}
真值 $T\in\{0,1\}$ 的先验概率为 $\Prob(T=1)=\pi\in(0,1)$。给定真值，各报告独立：以概率 $a$ 正确、以概率 $b$ 翻转、以概率 $1-a-b$ 缺失，其中 $a,b>0$、$a+b<1$。观察到 $h$ 个 $1$、$z$ 个 $0$ 和 $e$ 个缺失。求真值的后验赔率，并说明缺失是否提供信息。若缺失概率依赖真值，后一结论是否仍成立？

\exercise{PG1-4}{总和接近零的筛选}
$X,Z$ 独立标准正态，只保留满足 $|X+Z|<a$ 的观测，其中 $a>0$。求保留样本的相关系数，并求 $a\downarrow0$ 和 $a\to\infty$ 时的极限。

\exercise{PG2-1}{不等概率的第一次重复}
独立抽样有三种可能类型，概率分别为 $1/2,1/4,1/4$。令 $T$ 为第一次抽到已出现类型的时刻。求 $\E T$ 和 $\Prob(T=4)$。

\exercise{PG2-2}{圆环中的连续红段}
把 $m\ge1$ 张红卡和 $n\ge1$ 张黑卡均匀放在 $N=m+n$ 个标号圆环位置上，首尾相连。求红色段数的期望，以及全部红卡恰组成一个连续圆环段的概率。

\exercise{PG2-3}{同桶三元组的方差}
$m\ge3$ 个对象独立进入 $d$ 个桶，入桶概率为 $p_1,\ldots,p_d$，其中 $p_i\ge0$、$\sum p_i=1$。令 $C$ 为全部落入同一桶的无序三元对象组数。记 $s_r=\sum_i p_i^r$，求 $\E C$ 与 $\Var C$。

\exercise{PG2-4}{目标两侧的权重差}
把 $K\ge1$ 张目标卡与 $m\ge0$ 张普通卡均匀随机排列，普通卡 $i$ 带固定实数权重 $w_i$。令 $D$ 为首张目标卡之前的普通卡权重和，减去末张目标卡之后的普通卡权重和。求 $\E D$ 与 $\Var D$，允许权重为负。

\exercise{PG3-1}{增量相关的抽样成本}
$U_i$ 独立服从 $\operatorname{Unif}(0,1)$，在累计和首次超过 $1$ 时停止，次数为 $N$。第 $j$ 次抽样的成本是实际观测值 $U_j^2$。求总成本 $C=\sum_{j=1}^NU_j^2$ 的期望，并说明所用独立性为何不要求 $N$ 与增量独立。

\exercise{PG3-2}{含固定缩减的乘法停止}
$U_i$ 独立服从 $\operatorname{Unif}(0,1)$，令 $R_n=2^{-n}\prod_{i=1}^nU_i$、$R_0=1$。对 $0<x<1$，在首次 $R_n<x$ 时停止，次数为 $N$。给出 $N$ 的确定性上界、尾概率和均值的有限和表达式，并求 $x=1/4$ 时的完整分布。

\exercise{PG3-3}{区间内的更新次数}
$U_i$ 独立服从 $\operatorname{Unif}(0,1)$，$S_k=\sum_{i=1}^kU_i$。对 $0<a<b\le1$，令 $K$ 为落在 $(a,b]$ 内的正累计和 $S_k$ 的个数。求 $K$ 的完整分布和期望。

\pagebreak
\exercise{PG3-4}{剩余总量的最佳预测}
在题 \ref{prob:P24} 的模型中，已观察 $R_1,\ldots,R_n$，$n\ge0$，记 $Q_n=\sum_{j=1}^n(1-R_j)$。用这些观测预测完整总量 $Q$，求平方损失下的最优预测、条件预测误差方差，以及观察前的均方误差。说明为什么不能简单用部分和 $Q_n$ 作为最优预测。

\exercise{PG4-1}{两个均匀量的相对份额}
$U,V$ 独立服从 $\operatorname{Unif}(0,1)$。令 $W=U/(U+V)$，求 $W$ 的密度、均值和方差，并判断它是否均匀分布。

\exercise{PG4-2}{样本极值的中点}
$n\ge2$ 个独立样本均匀分布于 $(0,1)$，记最小值为 $L$、最大值为 $H$。求极值中点 $C=(L+H)/2$ 的密度、均值和方差。

\exercise{PG4-3}{随机点所在段的长度}
用 $n-1$ 个独立均匀切点将单位线段分成 $n\ge2$ 段，再独立均匀选取线段上的一个点。令 $L$ 为该点所在段的长度。求 $L$ 的密度、$\E L$ 与 $\E(1/L)$，并与从 $n$ 段中等概率选一段的平均长度比较。

\exercise{PG4-4}{三个方向的符号变化}
$X,Y$ 独立标准正态，$0<\theta<\pi/2$。分别取向量 $(X,Y)$ 在方向 $0,\theta,2\theta$ 上的投影符号，记为 $A,B,C\in\{-1,1\}$。令 $D_1=\ind\{A\ne B\}$、$D_2=\ind\{B\ne C\}$。求 $(D_1,D_2)$ 的联合分布、协方差及三个符号全相同的概率。
